% Generated by scripts/build_blueprint_tex.py from data/model.json -- do not edit by hand.
\begin{bpleanquote}{lean/Atlas/Architecture.lean}{31}{60}
\begin{bplean}[firstnumber=31]
\LCom{/-- The objects the endpoints talk about. `Point` stands for complex numbers `s`, `Char` for}
\LCom{Dirichlet characters of all positive moduli, `Zero χ s` for `L(s, χ) = 0`. Lean's `L`-functions}
\LCom{are total functions, so a value exists even at the principal pole. -/}
\LKw{structure} \LDecl{Plane} \LKw{where}
  Point : \LSort{Type}
  Char : \LSort{Type}
  Zero : Char → Point → \LSort{Prop}
  \LCom{/-- `Re s > 1/2`. -/}
  Right : Point → \LSort{Prop}
  \LCom{/-- `Re s = 1/2`. -/}
  Line : Point → \LSort{Prop}
  \LCom{/-- `Re s < 1/2`. -/}
  Left : Point → \LSort{Prop}
  \LCom{/-- `0 < Re s < 1`. -/}
  Strip : Point → \LSort{Prop}
  \LCom{/-- `(χ, s) = (1_N, 1)`: the principal character at the pole. -/}
  Pole : Char → Point → \LSort{Prop}
  Primitive : Char → \LSort{Prop}
  \LCom{/-- The primitive character inducing `χ`. -/}
  prim : Char → Char
  \LCom{/-- `χ ↦ χ⁻¹`. -/}
  inv : Char → Char
  \LCom{/-- `s ↦ 1 - s`. -/}
  refl : Point → Point
  \LCom{/-- The character of modulus one, whose `L`-function is `ζ`. -/}
  zeta : Char
  \LCom{/-- `s ∈ \{-2, -4, -6, …\}`. -/}
  Trivial : Point → \LSort{Prop}
  \LCom{/-- `s = 1`. -/}
  One : Point → \LSort{Prop}
\end{bplean}
\end{bpleanquote}
